% =====================================================================
%  1bat-ud1-14.tex  —  HORA 14: Repàs global i síntesi de la unitat (estacions)
% =====================================================================
\horatitol{Sessió 14 — Repàs global i síntesi de la unitat}%
{Repàs actiu per quatre estacions, una per bloc, i mapa conceptual que connecta cada eina amb el seu context real.}

\subsection{Idea clau}

Repassem \textbf{tota la unitat} de manera activa, rotant per \textbf{quatre
estacions}, una per cada bloc treballat. Abans, cada grup farà un \textbf{mapa
conceptual} connectant les eines matemàtiques amb els contextos reals on les hem
fet servir.

\begin{idea}
\textbf{Les quatre estacions}
\begin{enumerate}[leftmargin=1.6em, itemsep=2pt]
  \item \textbf{Intervals:} llegir un fragment real breu i escriure'n l'interval.
  \item \textbf{Inequacions:} resoldre'n una i comprovar-la per substitució.
  \item \textbf{Sistemes:} un cas amb dues condicions simultànies (intersecció).
  \item \textbf{Notació científica:} una conversió i una operació amb xifres grans.
\end{enumerate}
\textbf{El mapa de la unitat} connecta: nombres reals i recta $\rightarrow$ intervals
$\rightarrow$ inequacions $\rightarrow$ notació científica, amb els contextos (bo
social, frontera dels $18$, ajut al lloguer, pressupostos socials).
\end{idea}

\subsection{Les estacions}

\begin{exemple}[Estació 1 — intervals]
«Un ajut és per a infants \textbf{de $3$ a $12$ anys, inclosos}.» L'interval és
$[3,12]$ (tots dos extrems tancats). Si fos «menors de $12$», seria $[3,12)$.
\end{exemple}

\begin{exemple}[Estació 4 — notació científica]
Convertim $\num{4200000}$ a notació científica: $4,2\times 10^{6}$. I operem:
$\dfrac{4,2\times 10^{6}}{2\times 10^{3}}=2,1\times 10^{3}=\num{2100}$.
\end{exemple}

\subsection{Exercicis resolts}

\begin{exresolt}[model de l'Estació 2: inequació amb comprovació]
Resol $-3x+2\le 11$, dóna l'interval i comprova-la.
\solucio
\[
-3x\le 11-2 \;\Rightarrow\; -3x\le 9.
\]
Dividim per $-3$ (negatiu) i \textbf{girem}: $x\ge -3$, interval $[-3,+\infty)$.
\textbf{Comprovació} amb $x=0$: $-3\cdot 0+2=2\le 11$, cert.
\resposta{$x\ge -3$, interval $[-3,+\infty)$.}
\end{exresolt}

\begin{exresolt}[model de l'Estació 3: dues condicions]
Una beca és per a estudiants de $16$ a $21$ anys (inclosos) \emph{i} amb renda
familiar $\le\num{20000}\eur$. Decideix per a un estudiant de $18$ anys amb renda
$\num{21000}\eur$.
\solucio
Edat $18\in[16,21]$: compleix. Renda $\num{21000}\le\num{20000}$? \textbf{No}
($\num{21000}>\num{20000}$). Com que falla una condició, \textbf{no} té dret a la beca.
\resposta{No: compleix l'edat però no la renda ($\num{21000}>\num{20000}$).}
\end{exresolt}

\subsection{Exercicis per fer}

\noindent\textit{Una tasca per estació. Rota i resol-les totes.}

\begin{exercicis}
\begin{exlist}
  \item \textbf{(Estació 1)} Escriu l'interval de cada condició i digues com és cada
        extrem:
        \begin{enumerate}[label=\alph*), leftmargin=1.6em, topsep=2pt]
          \item «renda com a màxim $\num{15000}\eur$»
          \item «edat de $18$ a menys de $30$ anys»
        \end{enumerate}
  \item \textbf{(Estació 2)} Resol i comprova: $2x-5>3x-9$.
  \item \textbf{(Estació 2)} Resol i dóna l'interval: $4(x-1)\le 2x+6$.
  \item \textbf{(Estació 3)} Troba l'interval que compleix alhora $x\ge -2$ i $x<5$,
        i representa'l a la recta.
  \item \textbf{(Estació 4)} Converteix $\num{36000000}$ a notació científica i
        calcula $\dfrac{3,6\times 10^{7}}{9\times 10^{3}}$.
  \item \textbf{(Mapa de la unitat)} Completa la connexió de cada eina amb el seu
        context real:
        \begin{center}
        \renewcommand{\arraystretch}{1.4}
        \begin{tabular}{p{4.6cm} p{8cm}}
        \toprule
        \textbf{Eina matemàtica} & \textbf{Context real on l'hem fet servir} \\
        \midrule
        Intervals & \rule{7.4cm}{0.4pt} \\
        Inequacions & \rule{7.4cm}{0.4pt} \\
        Sistemes / interseccions & \rule{7.4cm}{0.4pt} \\
        Notació científica & \rule{7.4cm}{0.4pt} \\
        \bottomrule
        \end{tabular}
        \end{center}
\end{exlist}
\end{exercicis}

% ---------------------------------------------------------------------
\fullsolucions{Sessió 14}
\begin{solucions}
\begin{sollist}
  \item (a) $(-\infty,\,\num{15000}]$, extrem dret tancat. \quad
        (b) $[18,30)$, el $18$ tancat i el $30$ obert.
  \item $2x-3x>-9+5 \Rightarrow -x>-4$; girem: $x<4$, interval $(-\infty,4)$.
  \item $4x-4\le 2x+6 \Rightarrow 2x\le 10 \Rightarrow x\le 5$, interval $(-\infty,5]$.
  \item Intersecció: $-2\le x<5$, interval $[-2,5)$; punt ple a $-2$, buit a $5$.
  \item $\num{36000000}=3,6\times 10^{7}$; \quad
        $\dfrac{3,6\times 10^{7}}{9\times 10^{3}}=0,4\times 10^{4}=4\times 10^{3}=\num{4000}$.
  \item Resposta oberta. Idees: intervals $\rightarrow$ bo social / frontera dels $18$;
        inequacions $\rightarrow$ llindars de renda; sistemes $\rightarrow$ ajut al
        lloguer jove (edat i renda); notació científica $\rightarrow$ pressupostos
        socials i població.
\end{sollist}
\end{solucions}
